\subsection{Objetivo 4 --- construção dos pools}

\textit{Construir pools com algoritmos simples e amplamente conhecidos, como Random
Forest e Perceptron, induzindo diversidade por variações de treinamento.}

Cinco modos, todos com $M=100$ (protocolo completo na Seção~1):

\begin{description}
\item[PGDCS completo.] Pool de Perceptrons lineares gerado por algoritmo genético
que otimiza medidas de complexidade dos \emph{bags}, com a etapa de votação das
medidas (\texttt{get\_best\_types}) ativa --- a especificação original do método.
\item[PGDCS-F1/T1.] Idêntico, com as medidas fixas em F1 e T1 em vez de votadas por
base. É a variante já publicada em versões anteriores deste relatório.
\item[Bags aleatórios.] Controle: mesmos bags, mesmo Perceptron, sem busca genética.
\item[Bagging.] Árvores de decisão sobre reamostragens \emph{bootstrap}.
\item[Random Forest.] Árvores com subamostragem adicional de atributos.
\end{description}

Os cinco atendem ao objetivo (algoritmos simples; diversidade por variação de
subconjuntos de dados, de atributos e da própria seleção dos \emph{bags}) e cobrem
uma faixa larga de força do classificador-base, o que se mostrou determinante nas
seções seguintes.

A medida de diversidade merece destaque porque o número bruto inverte a conclusão:

\begin{table}[htbp]
\centering
\caption{Diversidade: o Double Fault bruto contra o índice normalizado.}
\begin{tabular}{lrrrr}
\toprule
Modo & Acc.\ individual & DF bruto & $e^{2}$ esperado & \dfe \\
\midrule
PGDCS completo & 0.7042 & 0.1571 & 0.0875 & 2.12 \\
PGDCS-F1/T1    & 0.7052 & 0.1566 & 0.0869 & 2.13 \\
Bags aleatórios& 0.7046 & 0.1565 & 0.0873 & \textbf{2.10} \\
Random Forest  & 0.7806 & 0.1069 & 0.0481 & 2.92 \\
Bagging        & 0.7931 & 0.1094 & 0.0428 & 3.98 \\
\bottomrule
\end{tabular}
\end{table}

Pelo DF bruto os três pools de Perceptron seriam os \emph{menos} diversos (valores
mais altos, $\approx0.156$--$0.157$); pelo índice normalizado eles são os
\emph{mais} diversos (valores mais baixos, $2.10$--$2.13$; Friedman
$p=0.00026$, $k=5$). O DF absoluto sozinho leva ao diagnóstico errado, porque
classificadores mais fracos erram mais e portanto erram juntos mais.

\paragraph{O confounder geração $\times$ classificador-base, resolvido.} A versão
anterior deste relatório não podia separar ``o GA gera diversidade útil'' de ``o
Perceptron é mais fraco individualmente'' --- os dois efeitos andavam juntos.
Comparando os cinco modos par a par (Wilcoxon pareado por base, $N=29$):

\begin{table}[htbp]
\centering\small
\caption{Contrastes entre modos de geração --- nenhuma métrica difere entre PGDCS
completo, PGDCS-F1/T1 e bags aleatórios; todas diferem entre bags aleatórios e
Bagging.}
\begin{tabular}{lrrr}
\toprule
Contraste & \On{1} & $\mvr$ & $\dfe$ \\
\midrule
PGDCS completo \emph{vs} F1/T1 & $p=0.39$ & $p=0.88$ & $p=0.29$ \\
F1/T1 \emph{vs} bags aleatórios & $p=0.59$ & $p=0.15$ & $p=0.13$ \\
Bags aleatórios \emph{vs} Bagging & $p=0.12$ & $\mathbf{p=0.0038}$ & $\mathbf{p=0.0004}$ \\
\bottomrule
\end{tabular}
\end{table}

\textbf{Nenhuma das duas variáveis de geração dentro do Perceptron --- votar as
medidas de complexidade ou buscar com o GA --- produz uma diferença estatisticamente
sustentada em \On{1}, $\mvr$ ou $\dfe$.} A diferença aparece inteira no contraste
seguinte, bags aleatórios \emph{vs} Bagging: mesma ausência de busca dos dois lados,
único fator restante é o classificador-base, e é aí que $\mvr$ cai $0.066$ e $\dfe$
cai de $2.10$ para $3.98$. \textbf{O eixo que organiza os resultados deste
relatório é o classificador-base, não o método de geração.} A tese de que a busca
do GA por diversidade produz especialistas que a votação desperdiça não se sustenta
nestes dados: o mesmo padrão (\On{1} altíssimo, $\mvr$ baixo, folga grande) aparece
igual com ou sem busca, e some quando o classificador-base muda para árvore.

Isso também fecha, por medição, se o atalho \texttt{F1/T1} custa resultado: não
custa. Reativar a votação completa do PGDCS gastou $10{,}7$h de relógio
concentradas em duas bases (Magic e WDVG) para produzir um pool
estatisticamente indistinguível do que as medidas fixas já entregavam. Curiosamente,
a votação raramente escolhe F1/T1: o par aparece em apenas $4$ dos $290$ folds
($1{,}4\%$); os pares mais votados são combinações com F3, F4, LSC e N-medidas
(Seção~8, objetivo~7). O atalho não é uma boa adivinhação da votação --- é uma
substituição que, apesar de escolher medidas diferentes, chega ao mesmo resultado.
